\subsection{Complements on compact spaces and inverse limits}

PROPOSITION 1. --- Let X and Y be two topological spaces and f a continuous mapping of X into Y. Let \((\mathrm{K}_{\alpha})_{\alpha \in \mathrm{A}}\) be a decreasing directed family of compact subsets of X, with intersection K. Then \(f(\mathrm{K}) = \bigcap f(\mathrm{K}_{\alpha})\) .

For, let \(y\) be a point of \(\bigcap_{\alpha \in \mathbf{A}} f(\mathbf{K}_{\alpha})\) ; for every \(\alpha \in \mathbf{A}\) , the set \(\mathbf{L}_{\alpha} = \mathbf{K}_{\alpha} \cap \overline{f}^{-1}(y)\) is compact and nonempty. The family \((\mathbf{L}_{\alpha})_{\alpha \in \mathbf{A}}\) is directed downward, therefore its intersection \(\mathbf{L}\) is nonempty. Now, \(\mathbf{L} = \mathbf{K} \cap \overline{f}^{-1}(y)\) , whence \(y \in f(\mathbf{K})\) . We have thus proved the inclusion \(f(\mathbf{K}) \supset \bigcap_{\alpha \in \mathbf{A}} f(\mathbf{K}_{\alpha})\) , and the reverse inclusion is obvious.

PROPOSITION 2. --- Let there be given an inverse system \((\mathrm{T}_{i}, p_{ij})\) of topological spaces indexed by I, a topological space T, and a coherent and separating family of continuous mappings \(p_{i}: T \to T_{i}\) . Then:

\begin{enumerate}
\def\labelenumi{\alph{enumi})}
\item
  For every compact subset K of T, one has \(\mathbf{K} = \bigcap_{i\in I}\overline{p_i}^{-1}\big(p_i(\mathbf{K})\big)\) .
\item
  Let K and L be two disjoint compact subsets of T. There exists an \(i \in I\) such that \(p_{j}(K)\) and \(p_{j}(L)\) are disjoint for \(j \geqslant i\) .
\item
  Let \(x\) be a point of \(\bigcap_{i \in I} \overline{p_i}^{-1}(p_i(K))\) ; for every \(i \in I\) , the set \(K_i\) of points \(y\) in \(K\) such that \(p_i(y) = p_i(x)\) is a nonempty closed subset of \(K\) . For \(i \leqslant j\) we have \(K_i \supset K_j\) , and, since \(K\) is compact, the set \(\bigcap_{i \in I} K_i\) is therefore nonempty. Let \(y\) be a point of \(\bigcap_{i \in I} K_i\) ; we have \(y \in K\) and \(p_i(y) = p_i(x)\) for all \(i \in I\) , whence \(y = x\) ; finally, \(x \in K\) , which proves the inclusion \(K \supset \bigcap_{i \in I} \overline{p_i}^{-1}(p_i(K))\) ; the reverse inclusion is obvious.
\item
  For every \(i \in \mathbf{I}\) , set \(\mathbf{M}_i = \overline{p_i}^{-1}(p_i(\mathbf{K})) \cap \mathbf{L}\) ; this is a closed subset of the compact space \(\mathbf{L}\) , we have \(\mathbf{M}_i \supset \mathbf{M}_j\) for \(i \leqslant j\) , and, by \(a\) ),
\end{enumerate}

\[
\bigcap_ {i \in I} M _ {i} = K \cap L = \varnothing .
\]

Consequently, there exists an index \(i\) such that \(\mathbf{M}_i = \varnothing\) . For \(j \geqslant i\) , we have \(\overline{p_j}^{-1} (p_j(\mathbf{K})) \cap \mathbf{L} = \mathbf{M}_j \subset \mathbf{M}_i = \varnothing\) , whence \(p_j(\mathbf{K}) \cap p_j(\mathbf{L}) = \varnothing\) .
% semantic-fragments: legacy-input-seam
\relax{}
